\documentclass[10pt,a4paper]{article}
\usepackage[utf8]{inputenc}
\usepackage[a4paper,left=3cm,right=2cm]{geometry}
\usepackage{amsmath}
\usepackage[thmmarks, amsmath, thref]{ntheorem}
\usepackage{amsfonts}
\usepackage{amssymb}
\usepackage{fancyhdr}
\pagestyle{fancy}
\usepackage{anysize}
\usepackage{graphicx}
\usepackage{float}
\usepackage{hyperref}
\usepackage{enumitem}
\usepackage{verbatim}
\newtheorem{theorem}{Theorem}
\theoremsymbol{\ensuremath{\blacksquare}}
\newtheorem*{solution}{Solución:}

\usepackage{xcolor}

% Margenes y formalidades

\textwidth 6.25in % Margen Derecho y ancho del texto
\textheight 8.5in
\oddsidemargin 0in % Margen Izquierdo
\headheight 0in % Largo del texto
\leftmargin 1in
\parindent 0em % Salto entre la indentación (Sangría)https://www.overleaf.com/4519127669sdpptkhvsbfn
\topmargin -0.5in % Margen Superior
\parskip 0.5ex % Salto entre parrafos
\baselineskip 1pt



\lhead{
\includegraphics[scale=0.08]{Figuras/logoUSM.png}} % texto izquierda de la cabecera
%\chead{}                                      % texto centro de la cabecera
\rhead{Universidad Técnica Federico Santa María\\Ingeniería Comercial } % número de página a la derecha
\renewcommand{\headrulewidth}{0.4pt} % grosor de la línea de la cabecera
\renewcommand{\footrulewidth}{0.4pt} % grosor de la línea del pie

\begin{document}
\vspace{5cm}



\begin{center}
\begin{large}
\textbf{\\CERTAMEN 1}
\end{large}\\
\end{center}

\begin{center}
\begin{large}
\textbf{Econometría ICS-294}
\end{large}\\
\vspace{0.2cm}
I Semestre 2025\\
Profesores: Jocelyn Tapia Stefanoni, Franco Antonucci Fernández y Ricardo Pruneda Urz\'ua\\
%Ayudante: 
\end{center}

\textbf{Instrucciones:} 

\begin{itemize}

    \item Usted dispone de 90 minutos  para intentar los 100 puntos del certamen.
    \item Cuide la presentación y redacción de sus preguntas.
    \item No puede utilizar celular, computador o tablet, pero sí una calculadora TI.
     
\end{itemize}

\section{RdA3.2  Preguntas cortas (20 puntos)}
{\bf Elija solo 2 preguntas y responda. }
\begin{enumerate}
    \item (10 puntos) Dos economistas, Milana y Maxi, están analizando datos de una muestra de personas de todos los países de América Latina para estudiar cómo influye la \textit{cantidad de aplicaciones abiertas en el celular} sobre la \textit{duración de la batería}. 
    Debido a posibles errores de medición en el proceso de recolección de datos, surge un debate entre ambas. Milana sostiene que, desde el punto de vista econométrico y en términos de obtener un mejor estimador, sería menos problemático que el error de medición estuviera en la variable \textit{duración de la batería}. Por su parte, Maxi argumenta que, por el contrario, sería menos grave que el error estuviera en la variable \textit{cantidad de aplicaciones abiertas en el celular}. ¿Quién tiene la razón? ¿Por qué? Puede apoyarse con herramientas matemáticas.

{\color{blue}
     Las economistas están estimando el siguiente modelo (1 punto por expresar el modelo entendiendo cuál es la independiente y la dependiente; no es necesario plantearlo matemáticamente, podía ser en palabras):

    $$Duracion_i = \beta_0 + \beta_1 \cdot Apps_i + u_i$$

    Donde
    \begin{itemize}
        \item \textit{Duración} = duración de la batería (variable dependiente)
        \item \textit{Apps} = cantidad de aplicaciones abiertas en el celular (variable independiente)
        \item \textit{u} = término de error (variables no observables)
    \end{itemize}

    Por lo tanto Milana indica que en es menos grave que el error de medición esté en la variable dependiente, mientras que Maxi indica que es menos grave cuando el error de medición está en la variable independiente. (2 puntos, entender qué plantea cada una) \\
    Cuando el error de medición se encuentra en la variable dependiente se obtendrán estimadores insesgados, pero menos eficientes, es decir tendrán mayor varianza. Por otro lado, cuando el error de medición se encuentra en la variable independiente, entonces se obtendrán estimadores sesgados. Puesto que si $X = X^* + u$, donde la variable de interés es $X^*$ pero sólo observamos $X$, entonces si el modelo verdadero viene dado por $Y = \beta_0 + \beta_1 \cdot X^* + e$ (1) se podrá estimar $Y = b_0 + b_1 \cdot X + v$. (2 puntos por entender qué genera cada error de medición)\\

    Luego, 
    $$E\left[\hat{b}_{1}\right]=\frac{\operatorname{Cov}(X, Y)}{\operatorname{Var}(X)}=\frac{\operatorname{Cov}\left(X^{*}+u, Y\right)}{\operatorname{Var}\left(X^{*}+u\right)}$$

    Dado que el error es clásico: $\operatorname{Cov}(u, Y)=0$ por lo que:
    $$E\left[\hat{b}_{1}\right]=\frac{\operatorname{Cov}\left(X^{*}, Y\right)}{\operatorname{Var}\left(X^{*}+u\right)}$$
    Además, el denominador lo podemos escribir como la suma de varianzas:
    $$\operatorname{Var}\left(X^{*}+u\right)=\operatorname{Var}\left(X^{*}\right)+\operatorname{Var}(u)$$
    
    Reemplazando $Y$ por el modelo original (1):
    $$E\left[\hat{b}_{1}\right]=\frac{\operatorname{Cov}\left(X^{*}, \beta_{0}+\beta_{1} X^{*}+\epsilon\right)}{\operatorname{Var}\left(X^{*}\right)+\operatorname{Var}(u)}$$
    separando las covarianzas:
    $$E\left[\hat{b}_{1}\right]=\frac{\beta_{1} \operatorname{Cov}\left(X^{*}, X^{*}\right)+\operatorname{Cov}\left(X^{*}, \epsilon\right)}{\operatorname{Var}\left(X^{*}\right)+\operatorname{Var}(u)}$$
    Además, sabemos que $\operatorname{Cov}\left(X^{*}, \epsilon\right)=0$. y $\operatorname{Cov}\left(X^{*}, X^{*}\right)=\operatorname{Var}\left(X^{*}\right)$.
    Entonces:
    $$E\left(\hat{b}_{1}\right)=\beta_{1} \frac{\operatorname{Var}\left(X^{*}\right)}{\operatorname{Var}\left(X^{*}\right)+\operatorname{Var}(u)}=a \cdot \beta_{1}$$
    Por lo tanto, al ser insesgado el estimador de $\beta_1$ cuando el error de medición está en la variable dependiente, y al ser sesgado cuando está en la independiente, entonces Milana tiene la razón. \\(5 puntos por concluir bien)\\
    
    (No es necesario hacer toda la matemática para obtener el puntaje, si la explicación está bien)
    %(causalidad- error de medicion-var. omitida)
    
}

    %intervalos de confianza
    \item (10 puntos) Describa el esquema de muestreo utilizado por Miguel y Kremer (2004) en su estudio {\it Worms: Identifying Impacts on Education and Health in the Presence of Treatment Externalities}.
    ¿Qué ventajas tiene este diseño para la identificación causal de los efectos del tratamiento?

    {\color{blue} Miguel y Kremer (2004) utilizaron un diseño de muestreo por conglomerados (cluster sampling), donde la unidad de aleatorización fue la escuela para capturar externalidades (4 puntos). Trabajaron con 75 escuelas primarias rurales en Kenia, divididas aleatoriamente en tres grupos: dos grupos recibieron tratamiento en distintos años (1998 y 1999), y un tercer grupo actuó como control durante el estudio.
    La aleatorización por escuelas reduce el riesgo de sesgo de selección y permite atribuir diferencias causales en resultados al tratamiento.(6 puntos)}

\item  (10 puntos) Josefina se encuentra trabajando para una importante constructora nacional en el \'area de finanzas, el jefe de
proveedores le comenta que la media de cuentas por cobrar es de 220.000.000 de pesos. Josefina tiene una  muestra aleatoria de 500 facturas con media igual a 195.000.000 de pesos y una desviaci\'on
est\'andar de 1.000.000. Con un 90\% de confianza ¿es correcta la afirmacion del jefe?


    $$\color{blue}{\bar{x} \pm z_{\frac{\alpha}{2}}\frac{\sigma}{\sqrt{n}}}$$

{\color{blue}$[194.926.209,8 ; 195.073.790,2]$ (6 puntos)}


   {\color{blue}  El valor que comenta el jefe es 220.000.000 el cual no está en el intervalo de confianza calculado, así que  90\% de confianza la afirmación es incorrecta. (4 puntos)}

\item (10 puntos) La  suma de cuadrados de los residuos es:
    \[
    SSR = \sum_{i=1}^{n} \hat{u}_i^2 = \hat{u}^T \hat{u}
    \]
    
Demuestre que esta puede reescribirse como:
    \[
    SSR = y^T y - \hat{y}^T y
    \]
    donde \( \hat{y} = X \hat{\beta} \) representa los valores ajustados del modelo, y \( \hat{u} = y - \hat{y} \) es el vector de residuos.

{\color{blue}  \[
\hat{u}^T \hat{u} = y^T y - 2 \hat{y}^T y + \hat{y}^T \hat{y}.
\]
(2 puntos)
Ahora notamos que \( \hat{y} = X \hat{\beta} \) y \( \hat{y}^T \hat{y} = \hat{\beta}^T X^T X \hat{\beta} \), mientras que \( \hat{y}^T y = \hat{\beta}^T X^T y \). Como \( \hat{\beta} = (X^T X)^{-1} X^T y \), entonces:

\[
\hat{y}^T y = \hat{\beta}^T X^T y = y^T X (X^T X)^{-1} X^T y = \hat{y}^T \hat{y}.
\]

(4 puntos)
Por lo tanto:

\[
\hat{y}^T \hat{y} = \hat{y}^T y.
\]


Sustituyendo esto en la expresión del SSR:

\[
\hat{u}^T \hat{u} = y^T y - 2 \hat{y}^T y + \hat{y}^T y = y^T y - \hat{y}^T y.
\]

(4 puntos)
}
\end{enumerate}

\section{RdA3.2 Ejercicios (80 puntos)}
{\bf Elija 4 de las siguientes preguntas y responda.} 
\begin{enumerate}

  \item La gerencia de Recursos Humanos de una reconocida empresa de servicios está interesada en mejorar la calidad de la atención mediante la introducción de una nueva herramienta de trabajo. Para ello, ha invitado a los trabajadores con más de 2 años en la empresa a inscribirse voluntariamente en una capacitación. De un total de 15 empleados, 5 decidieron capacitarse. Un año después, el salario promedio del grupo capacitado asciende a US\$1.850, mientras que el del grupo que no se capacitó es de US\$1.600.
\begin{enumerate}
      \item  (5 puntos) ¿Podemos concluir que existe una relación \textit{causal} entre participar en la capacitación y el aumento del salario? Justifique.

{\color{blue}  No se puede concluir una relación causal (2 puntos), dado que el grupo de tratamiento no es aleatorio y puede depender de del trabajador, como motivación, inteligencia, entre otros. (3 puntos)}
    \item  (7 puntos) Halle el valor del coeficiente $\hat{\beta}$ en un modelo de regresión lineal simple que explique el salario en función de la variable dicotómica ``capacitación'' (1 = capacitado, 0 = no capacitado).

{\color{blue}    $$\hat{\beta} = 1850-1600 = 250$$}
    \item  (8 puntos) Se sabe que las \textit{desviaciones estándar} del salario para el grupo capacitado y para el no capacitado  son  US\$250  y  US\$200 respectivamente. Con un nivel de confianza del 95\%, ¿es estadísticamente significativa la diferencia de medias entre ambos grupos?

    {\color{blue}  El intervalo de confianza esta dado por:
    $$IC = (\mu_T - \mu_{NT}) \pm Z_{\alpha/2}\sqrt{\frac{\sigma_T^2}{n_T}+\frac{\sigma^2_{NT}}{n_{NT}}}$$
    $$IC = 1850-1600\pm 1.96\sqrt{\frac{250^2}{5}+\frac{200^2}{10}}$$

    (5 puntos)

    $$IC = [ -1.767 ;\quad 501.767]$$
    La diferencia no es estadisticamente significativa ya que contiene el valor 0 en el intervalo de confianza.  (3 puntos)}
\end{enumerate}
    

\item (20 puntos) Considere una muestra de tamaño \( n \), pero en la cual las variables observadas presentan errores de medición:

\[
Y_i^* = Y_i + u_i, \quad \text{donde } u_i \sim \mathcal{N}(0,\,1), \quad \text{y}
\]
\[
X_i^* = X_i + c, \quad \text{donde } c \neq 0 \text{ es una constante}
\]

El modelo verdadero (sin errores de medición) es:
\[
Y_i = \beta_0 + \beta_1 X_i + u_i
\]

Compare los estimadores de Mínimos Cuadrados Ordinarios (MCO) obtenidos utilizando los datos observados \((Y_i^*, X_i^*)\) versus aquellos que se obtendrían si se observaran las variables verdaderas \((Y_i, X_i)\), en una muestra de tamaño \( n \). Analice el efecto que tienen los errores de medición en la estimación de los parámetros del modelo.

{\color{blue}
$$\hat{\beta_1^*} = \frac{\sum (X_i^*-\Bar{X^*})(Y_i^*-\Bar{Y^*})}{\sum (X_i^*-\Bar{X^*})^2}$$
$$\hat{\beta_1^*} = \frac{\sum (X_i+c-(\Bar{X}+c))(Y_i+u-(\Bar{Y}+\Bar{u}))}{\sum (X_i+c-(\Bar{X}+c))^2}$$

(4 puntos)
$$\hat{\beta_1^*} = \frac{\sum (X_i-\Bar{X})(Y_i+u-(\Bar{Y}+\Bar{u}))}{\sum (X_i-\Bar{X})^2}$$
$$\hat{\beta_1^*} = \frac{\sum (X_i-\Bar{X})(Y_i-\Bar{Y})}{\sum (X_i-\Bar{X})^2} + \frac{\sum (X_i-\Bar{X})(u-\Bar{u})}{\sum (X_i-\Bar{X})^2}$$
$$\hat{\beta_1^*} = \hat{\beta_1} + \frac{\sum (X_i-\Bar{X})(u-\Bar{u})}{\sum (X_i-\Bar{X})^2}$$

(4 puntos)

El t\'ermino $\frac{\sum (X_i-\Bar{X})(u-\Bar{u})}{\sum (X_i-\Bar{X})^2}$ es igual a 0 en terminos esperados, por lo que:
 $$E(\hat{\beta_1^*}) = E(\hat{\beta_1})$$

 (2 puntos) No era necesario poner el insesgamiento.
}


\item Considere el siguiente modelo de regresión poblacional:

\[
Y_i = \beta_0 + \beta_1 X_i + \varepsilon_i
\]

\begin{enumerate}
    \item[(a)] (10 puntos) Suponiendo que se dispone de una muestra aleatoria de tamaño \( n \), demuestre que los estimadores  MCO para los parámetros \( \beta_0 \) y \( \beta_1 \) están dados por:
    \[
    \hat{\beta}_1 = \frac{\sum_{i=1}^{n}(X_i - \bar{X})(Y_i - \bar{Y})}{\sum_{i=1}^{n}(X_i - \bar{X})^2}, \qquad
    \hat{\beta}_0 = \bar{Y} - \hat{\beta}_1 \bar{X}
    \]

\color{blue} Dado el modelo lineal simple:
\[
y_i = \beta_0 + \beta_1 x_i + \varepsilon_i, \quad i = 1, 2, \ldots, n
\]
Queremos encontrar los estimadores \(\hat{\beta}_0\) y \(\hat{\beta}_1\) que minimizan la suma de los errores al cuadrado:
\[
S = \sum_{i=1}^{n} (y_i - \beta_0 - \beta_1 x_i)^2
\]
\subsection*{Paso 1: Derivadas parciales}
Tomamos derivadas parciales de \(S\) con respecto a \(\beta_0\) y \(\beta_1\) y las igualamos a cero:
\[
\frac{\partial S}{\partial \beta_0} = -2 \sum_{i=1}^{n} (y_i - \beta_0 - \beta_1 x_i) = 0
\]
\[
\frac{\partial S}{\partial \beta_1} = -2 \sum_{i=1}^{n} x_i (y_i - \beta_0 - \beta_1 x_i) = 0
\]
\subsection*{Paso 2: Ecuaciones normales}
De la primera ecuación obtenemos:

(5 puntos)

\[
\sum y_i - n\beta_0 - \beta_1 \sum x_i = 0 \Rightarrow \beta_0 = \bar{y} - \beta_1 \bar{x}
\]
Donde:
\[
\bar{x} = \frac{1}{n} \sum x_i, \quad \bar{y} = \frac{1}{n} \sum y_i
\]
Sustituimos en la segunda ecuación:
\[
\sum x_i (y_i - (\bar{y} - \beta_1 \bar{x}) - \beta_1 x_i) = 0
\Rightarrow \sum x_i (y_i - \bar{y} + \beta_1 \bar{x} - \beta_1 x_i) = 0
\]
Reordenando:
\[
\sum x_i(y_i - \bar{y}) + \beta_1 \sum x_i (\bar{x} - x_i) = 0
\]
Observamos que \(\sum x_i (\bar{x} - x_i) = \bar{x} \sum x_i - \sum x_i^2 = n \bar{x}^2 - \sum x_i^2 \neq 0\), entonces:
\[
\hat{\beta}_1 = \frac{\sum_{i=1}^{n} (x_i - \bar{x})(y_i - \bar{y})}{\sum_{i=1}^{n} (x_i - \bar{x})^2}
\]
(5 puntos)
\color{black}
    \item[(b)] (10 puntos) Indique y explique los supuestos necesarios para que los estimadores MCO sean insesgados. Luego, bajo dichos supuestos, pruebe que los estimadores de \( \beta_0 \) y \( \beta_1 \) son efectivamente insesgados.

\color{blue}
\begin{enumerate}
    \item Linealidad en parámetros
    \item Muestra aleatoria
    \item Condición de identificación
    \item Ortogonalidad
\end{enumerate}
(5 puntos)

\color{blue}
Teorema: Bajo los supuestos 1 a 4, el estimador MCO de $\beta$ es insesgado.
Demostración:
    \begin{align*}
        \hat{\beta} &= (X^{\top}X)^{-1}X^{\top}y \quad \text{(3)} \\
        &= (X^{\top}X)^{-1}X^{\top}(X\beta + u) \quad \text{(4)} \\
        &= \beta + (X^{\top}X)^{-1}X^{\top}u \quad \text{(5)} \\
        E(\hat{\beta}|X) &= \beta + (X^{\top}X)^{-1}X^{\top}E(u|X) \quad \text{(6)} \\
        &= \beta \quad \text{(7)}
    \end{align*}
(5 puntos)
Esto implica que el promedio de los estimadores $\hat{\beta}$ obtenidos con todas las muestras aleatorias posibles es igual a $\beta$.

    
\end{enumerate}

\vspace{0.5cm}

    \item Considere el siguiente modelo de regresión lineal con los supuestos estándares:

    \[
    Y_i = \beta_0 + \beta_1 X_{1i} + \beta_2 X_{2i} + u_i
    \]
    
    Se dispone de la siguiente información:
    
    \[
    X'X = \begin{pmatrix}
    33 & 0 & 0 \\
    0 & 40 & 20 \\
    0 & 20 & 60
    \end{pmatrix}, \quad
    X'Y = \begin{pmatrix}
    132 \\
    24 \\
    92
    \end{pmatrix}
    \]
    
    \[
    \sum_{i=1}^{N} Y_i^2 = 678, \quad \sum_{i=1}^{N} (Y_i - \bar{Y})^2 = 150
    \]

    \begin{enumerate}
        \item (3 puntos) ¿Cuál es el tamaño de la muestra?\\

        {\color{blue} 33, el valor $a_{11}$ de la matriz $X'X$.} (3 puntos, respuesta directa)\\
        
        \item (9 puntos) Calcule los estimadores de MCO del modelo anterior.\\
        \textbf{Hint:} \( Z = \begin{pmatrix} A &{\bf 0} \\ {\bf 0} & B \end{pmatrix} \) con \( A \) y \( B \) matrices cuadradas, entonces:\[Z^{-1} = \begin{pmatrix} A^{-1} & {\bf 0} \\{\bf 0} & B^{-1} \end{pmatrix}\]

        {\color{blue} Sabemos que los estimadores de mínimos cuadrados ordinarios (MCO) están dados por:
    
        \[
        \hat{\beta} = (X'X)^{-1} X'Y
        \]
    
        Calculando los estimadores (2 puntos por plantear):
    
        \[
        \hat{\beta} = \begin{pmatrix}
        33 & 0 & 0 \\
        0 & 40 & 20 \\
        0 & 20 & 60
        \end{pmatrix}^{-1}
        \begin{pmatrix}
        132 \\
        24 \\
        92
        \end{pmatrix}
        \]

        (2 puntos por calcular la inversa)
        \[
        \hat{\beta} = \begin{pmatrix}
        \frac{1}{33} & 0 & 0 \\
        0 & 0.03 & -0.01 \\
        0 & -0.01 & 0.02
        \end{pmatrix}
        \begin{pmatrix}
        132 \\
        24 \\
        92
        \end{pmatrix}
        \]

        (2 puntos por llegar al vector)
        \[
        \hat{\beta} = \begin{pmatrix}
        4 \\
        -0.2 \\
        1.6
        \end{pmatrix}
        \]
        
        Así que los estimadores son (1 punto por cada estimador indicado):
        
        \[
        \hat{\beta_0} = 4, \quad \hat{\beta_1} = -0.2, \quad \hat{\beta_2} = 1.6
        \]}
        \item (4 puntos) Indique la ecuación del modelo predicho e interprete.

        {\color{blue} La ecuación del modelo predicho es (2 puntos)

        $$\hat{y} = 4 - 0.2 \cdot x_1 + 1.6 \cdot x_2$$

        Como el modelo es lineal, esto indica que un aumento en una unidad de la variable $x_1$ hará caer a la variable $y$ en 0,2 unidades; mientras que si la variable $x_2$ aumenta en una unidad, la variable dependiente se incrementará en 1,6 unidades.} (2 puntos por la interpretación)\\
        
        \item (4 puntos) ¿Qué debería suceder con $R^2$ si agregamos más variables independientes a la medición?
        
        {\color{blue} El \( R^2 \) se puede escribir matemáticamente como:

        \[
        R^2 = 1 - \frac{\sum_{i=1}^{N} \hat{u_i}^2}{\sum_{i=1}^{N} (Y_i - \bar{Y})^2}
        \]

        Por lo tanto, el único valor que cambia cuando se agregan variables es \( \sum_{i=1}^{N} \hat{u_i}^2 \), que representa la función valor de la función objetivo, o sea, esta última evaluada en los parámetros que la minimizan. A su vez, por propiedad general de la optimización, la función objetivo {\bf no puede empeorar} si relajamos una restricción. Por ello, la suma de los residuos al cuadrado (la función objetivo de MCO) no puede empeorar si agregamos un regresor, ya que agregar un regresor es equivalente a relajar la restricción, en el sentido de que ahora el parámetro que acompaña a ese regresor no se impone igual a cero. Finalmente, al observar la expresión inicial, una suma de residuos al cuadrado más pequeña implica un mayor \( R^2 \). (4 puntos por interpretación)\\
}    
\end{enumerate}

\item  Supongamos que eres un investigador interesado en estudiar la relación entre la cantidad de horas de estudio (variable independiente X) y el puntaje obtenido en un examen (variable dependiente Y).
Tienes los siguientes datos de 6 estudiantes:

\begin{table}[h!]
\centering
\begin{tabular}{|c|c|}
\hline
\textbf{Horas de estudio (X)} & \textbf{Puntaje en el examen (Y)} \\
\hline
2 & 50 \\
3 & 60 \\
4 & 65 \\
5 & 70 \\
6 & 75 \\
7 & 76 \\
\hline
\end{tabular}
\label{tabla:estudio_puntaje}
\end{table}

\begin{enumerate}
    \item  (5 puntos) Use el método de los mínimos cuadrados para ajustar una línea recta $\hat{Y}=\hat{\beta_0}+\hat{\beta}_1X$.

    {\color{blue} \[
\bar{X} = 4.5, \quad \bar{Y} = 66.0
\]

\begin{center}
\begin{tabular}{|c|c|c|c|c|c|}
\hline
\textbf{X} & \textbf{Y} & \(X_i - \bar{X}\) & \(Y_i - \bar{Y}\) & \((X_i - \bar{X})(Y_i - \bar{Y})\) & \((X_i - \bar{X})^2\) \\
\hline
2 & 50 & -2.5 & -16.0 & 40.00 & 6.25 \\
3 & 60 & -1.5 & -6.0  & 9.00  & 2.25 \\
4 & 65 & -0.5 & -1.0  & 0.50  & 0.25 \\
5 & 70 & 0.5  & 4.0   & 2.00  & 0.25 \\
6 & 75 & 1.5  & 9.0   & 13.50 & 2.25 \\
7 & 76 & 2.5  & 10.0  & 25.00 & 6.25 \\
\hline
\textbf{Suma} & & & & \textbf{90.00} & \textbf{17.50} \\
\hline
\end{tabular}
\end{center}


\[
\hat{\beta}_1 = \frac{90}{17.5} = 5.14, \quad \hat{\beta}_0 = 66.0 - 5.14 \cdot 4.5 = 42.87
\]

El modelo estimado por mínimos cuadrados es:

\[
\hat{Y} = 42.87 + 5.14X
\]
}
   \item (5 puntos) ¿Cu\'al es el puntaje esperado en el examen para un estudiante que estudia 8 horas?

   {\color{blue} Utilizando el modelo estimado:

\[
\hat{Y} = 42.87 + 5.14 \cdot X,
\]

si un estudiante estudia 8 horas, el puntaje esperado en el examen es:

\[
\hat{Y} = 42.87 + 5.14 \cdot 8 = 42.87 + 41.12 = 83.99.
\]}
\item (5 puntos) Si un estudiante aumenta sus horas de estudios en 2 horas ¿cu\'anto se espera que aumente su puntaje?

{\color{blue} El coeficiente estimado de la variable \(X\) en el modelo es:

\[
\hat{\beta}_1 = {\Delta \hat{Y}\over \Delta X} =5.14
\]

Esto significa que, en promedio, por cada hora adicional de estudio, el puntaje en el examen aumenta en 5.14 puntos.

Por lo tanto, si un estudiante aumenta sus horas de estudio en 2 horas, se espera que su puntaje aumente en:

\[
\Delta \hat{Y} = 5.14 \cdot 2 = 10.28
\]
}

\item (5 puntos) Calcule el $R^2$ e interprete.

{\color{blue} 


\textbf{Calculamos  \( SST \)}

\[
SST = (-16)^2 + (-6)^2 + (-1)^2 + 4^2 + 9^2 + 10^2 = 256 + 36 + 1 + 16 + 81 + 100 = 490
\]

\textbf{ Calcular los valores ajustados \( \hat{Y}_i = 42.87 + 5.14 X_i \)}

\begin{itemize}
    \item \( X = 2 \Rightarrow \hat{Y} = 53.15 \Rightarrow (Y - \hat{Y})^2 = (50 - 53.15)^2 = 9.92 \)
    \item \( X = 3 \Rightarrow \hat{Y} = 58.29 \Rightarrow (60 - 58.29)^2 = 2.92 \)
    \item \( X = 4 \Rightarrow \hat{Y} = 63.43 \Rightarrow (65 - 63.43)^2 = 2.46 \)
    \item \( X = 5 \Rightarrow \hat{Y} = 68.57 \Rightarrow (70 - 68.57)^2 = 2.04 \)
    \item \( X = 6 \Rightarrow \hat{Y} = 73.71 \Rightarrow (75 - 73.71)^2 = 1.67 \)
    \item \( X = 7 \Rightarrow \hat{Y} = 78.85 \Rightarrow (76 - 78.85)^2 = 8.12 \)
\end{itemize}

\[
SSR = 9.92 + 2.92 + 2.46 + 2.04 + 1.67 + 8.12 = 27.13
\]


\[
R^2 = 1 - \frac{SSR}{SST} = 1 - \frac{27.13}{490} \approx 1 - 0.0554 = 0.9446
\]

 Aproximadamente el \textbf{94.5\%} de la variabilidad en el puntaje del examen se explica por el número de horas de estudio. 

}
\end{enumerate} 



    
\section{Formulario}
\begin{itemize}

\item El intervalo de confianza de $100(1-\alpha)\%$ para la media poblacional $\mu$ cuando $\sigma$ es \emph{\textbf{conocido}}:
$\color{red}{\bar{x} \pm z_{\frac{\alpha}{2}}\frac{\sigma}{\sqrt{n}}}$

\item El intervalo de confianza de $100(1-\alpha)\%$ para la diferencia de medias
$\mu_x-\mu_y$ es:
${\color{red}(\bar{X}-\bar{Y}) \pm z_{\frac{\alpha}{2}}\sqrt{\frac{\sigma_x^2}{n_x}+\frac{\sigma_y^2}{n_y}}}$

\item  Estimadores MCO.
{\color{red}
$$\hat{\beta}_0 = \bar{y} - \hat{\beta}_1 \bar{x} \quad \quad   \hat{\beta}_1 = \frac{\sum_{i=1}^{n}(x_i - \bar{x})(y_i-\bar{y})}{\sum_{i=1}^{n}(x_i - \bar{x})^2}$$}
    
\item 
	{\color{red}	\[ SST = \sum_{i=1}^{n} (y_i - \bar{y})^2 \quad \quad  SSE = \sum_{i=1}^{n} (\hat{y}_i - \bar{y})^2 \quad \quad SSR = \sum_{i=1}^{n} \hat{u}_i^2  \]}
		
		
    \item Vector de estimadores:    
{\color{red}    $\hat{\beta}=(X^TX)^{-1}X^Ty$}

\item {\color{red}\[
A = \begin{bmatrix}
a & b \\
c & d
\end{bmatrix}
\qquad
A^{-1} = \frac{1}{det(A)} \begin{bmatrix}
d & -b \\
-c & a
\end{bmatrix}
\]}
\end{itemize}

\begin{figure}[H]
    \centering
    \includegraphics[width=1.05\linewidth]{normal.png}
\end{figure}



\end{document}